\chapter{Reproducing the Results}\label{app:repro}

Everything this thesis reports can be regenerated from the repository
\url{https://github.com/besco1998/fanet-authbc}. This appendix says what is needed for each
part and what each command produces. The design rule behind it is stated in
\Cref{ch:implementation}: work that depends on a machine --- timings, simulations, radio and
energy measurements --- is run once and its raw output committed; everything derived from that
output is recomputed and compared on every run.

\section{Four paths, by cost}

{\small
\noindent\begin{tabularx}{\textwidth}{@{}l Y Y l@{}}
\toprule
path & reproduces & needs & time \\
\midrule
A, analytical & every byte, loss, exclusion, energy-model and optimiser result; every table and
figure, from the committed raw data & Python 3.12 or later, about 1\,GB of disk & minutes \\
B, 802.11 simulation & the channel-model validation and the capacity search &
ns-3.48 built from source, about 8\,GB of disk & hours \\
C, low-rate simulation & the LoRaWAN capacity and mobility runs & the ns-3 LoRaWAN module &
tens of minutes \\
D, hardware & primitive timings, powers, sender energy, two-radio airtime and loss &
two Raspberry~Pi~4, an INA219 current sensor, a host to read it & half a day \\
\bottomrule
\end{tabularx}
}

\noindent Path A alone reproduces every number in the results chapters, because the outputs
of paths B to D are committed and the analytical layer reads them.

\section{Path A}

\begin{verbatim}
git clone https://github.com/besco1998/fanet-authbc.git
cd fanet-authbc
make setup     # a virtual environment with the pinned dependencies
make all       # lint, type check, the test suite, the reproduction gate
make paper thesis
\end{verbatim}

\texttt{make all} ends with the reproduction gate: every derived result file is recomputed from
the current code and the committed raw data, and a single differing row fails the run.
\texttt{make paper} and \texttt{make thesis} first run \texttt{analysis/paper\_numbers.py},
which writes every number the two documents print into \texttt{numbers.tex} and stops if a
result file no longer supports a sentence that states a verdict. A number is changed by
re-running its experiment, never by editing the text.

\section{What is where}

{\small
\noindent\begin{tabularx}{\textwidth}{@{}l Y@{}}
\toprule
\texttt{src/authbc/} & the library: encodings, signature schemes, ledger, the two frame formats
with sender and receiver, the analytical models, the two simulators written for this work \\
\texttt{analysis/} & one script per group of results; \texttt{paper\_numbers.py}; figure scripts \\
\texttt{ns3/} & simulation scenarios (C++) and the drivers that run them over seeds \\
\texttt{hw/} & what runs on the boards: timing, the energy harness, the capture and reduction
of current samples, the two-radio link test \\
\texttt{experiments/} & configurations and run plans, including the lists of seeds \\
\texttt{results/raw/}, \texttt{results/hw/} & raw output, one file per experiment, each with a
header that records the environment and a hash of the configuration \\
\texttt{results/figures/} & figures, regenerated from the raw output \\
\texttt{tests/} & unit tests, published test vectors, the reproduction gate, and tests that hold
the text of the paper and of this thesis to the result files \\
\texttt{docs/} & specifications, the findings register, the expectation files of
\Cref{app:prereg}, the logbook \\
\bottomrule
\end{tabularx}
}

\section{Paths B to D}

The simulation scenarios are single source files copied into the simulator's \texttt{scratch}
directory; each driver names the scenario it runs, the seeds, and the file it writes. A run is
identified by its configuration and seed, and a stored run can be repeated bit for bit, which
is how a rebuilt simulator is checked before new runs are trusted. The hardware procedures ---
wiring, sampling, the order of a metered run and what makes a repetition unusable --- are in
\texttt{hw/}, written to be followed without the author.

\section{Software}

Python 3.12; \texttt{cbor2} pinned at 5.8.0, because a later release changed canonical output;
\texttt{cryptography} for Ed25519 and ECDSA and \texttt{blspy} for BLS12-381, each checked
against published test vectors before any timing; ns-3.48 with the \texttt{signetlabdei/lorawan}
module; PX4 v1.17.0 for the simulated flight of \Cref{ch:bytes}. The measured platform is a
Raspberry~Pi~4 (Cortex-A72) under a 64-bit Linux with the \texttt{performance} governor.
